\chapter*{Introduction générale}
\addcontentsline{toc}{chapter}{Introduction générale}
\markboth{Introduction générale}{}

\section*{Contexte}

L'évolution des systèmes d'information vers des architectures réparties s'accélère depuis le début des années deux mille. Les organisations développent leurs services autour de composants autonomes, interconnectés par des protocoles standardisés, plutôt qu'à l'intérieur d'applications monolithiques. Cette transformation répond à plusieurs impératifs : la mise à l'échelle horizontale, l'indépendance des cycles de livraison, la diversité technologique entre équipes, et l'interopérabilité entre organisations.

Dans cet écosystème, les services Web occupent une place centrale. Deux familles de protocoles dominent encore aujourd'hui les échanges machine à machine. La première, fondée sur le standard SOAP, privilégie un contrat strict décrit par un document WSDL et un format d'échange XML formel. La seconde, inspirée de l'approche REST formalisée par Roy Fielding dans sa thèse de doctorat \citep{fielding2000}, mise sur les ressources, les verbes du protocole HTTP et la négociation de contenu pour atteindre une grande flexibilité. Comprendre les forces respectives de ces deux familles, et savoir les mettre en oeuvre concrètement, constitue une compétence centrale pour tout ingénieur travaillant sur des systèmes distribués.

\section*{Problématique}

L'enseignement des services Web se limite souvent à une présentation théorique des standards, sans confrontation pratique des paradigmes sur un même cas métier. Or, c'est précisément cette confrontation qui révèle les compromis réels : verbosité des enveloppes XML, couplage induit par les contrats WSDL, simplicité de l'intégration HTTP, qualité des outils de test, gestion des erreurs, comportement face à l'absence d'un service distant. La question centrale de notre travail s'énonce ainsi : \emph{comment comparer rigoureusement les protocoles SOAP et REST en développant une même application métier dans ces deux paradigmes, tout en mettant en évidence les enjeux de composition de services et d'interopérabilité multi-langages ?}

\section*{Objectifs}

Pour répondre à cette interrogation, nous nous fixons quatre objectifs pédagogiques précis.

D'abord, mettre en oeuvre les deux protocoles sur un même scénario fonctionnel, en l'occurrence une chatroom multi-utilisateurs avec authentification, salons et messages. Cette identité fonctionnelle est la condition d'une comparaison valable.

Ensuite, composer deux services autonomes par phase qui communiquent entre eux \emph{via le protocole de la phase}, et non à travers une base de données partagée. Cette contrainte garantit que la composition est observable sur le réseau et que la cohérence des données passe par le protocole, non par une intégration au niveau du stockage.

Par ailleurs, démontrer une interopérabilité réelle en faisant cohabiter trois langages : Java, Node.js et Python. Une telle polyglossie reproduit les conditions de développement industriel où des équipes hétérogènes interviennent sur des services interdépendants.

Enfin, appliquer côté frontend le pattern \emph{Port et Adaptateurs} (architecture hexagonale) décrit par Alistair Cockburn \citep{cockburn2005}. Les pages applicatives doivent ignorer le protocole sous-jacent et basculer entre SOAP et REST par simple changement d'une variable d'environnement.

\section*{Méthodologie}

Notre démarche s'articule en deux phases iso-fonctionnelles. La première phase implémente l'application en SOAP, avec un service d'authentification écrit en Java JAX-WS et un service de chat écrit en Node.js. La seconde reproduit à l'identique les mêmes opérations en REST, en remplaçant les implémentations par Java JAX-RS et Python Flask. Les deux phases partagent le même frontend PHP, démontrant ainsi en pratique le pattern hexagonal.

Chaque phase fait l'objet d'une série de tests systématiques. Pour REST, nous utilisons l'outil Postman avec des scripts de tests réutilisables. Pour SOAP, nous mobilisons SoapUI avec des suites de tests chaînées par transferts de propriétés. Ces tests valident non seulement les opérations individuelles mais aussi la composition inter-service.

\section*{Plan du mémoire}

Le présent mémoire s'organise en six chapitres. Le premier chapitre expose le cadre théorique et l'état de l'art des deux paradigmes étudiés. Le deuxième présente l'analyse et la conception, avec une attention particulière portée au modèle conceptuel de données et au mécanisme de liaison cohérente entre les deux bases d'un même backend. Le troisième chapitre détaille la mise en oeuvre concrète, service par service. Le quatrième rapporte les campagnes de tests et illustre par des captures les échanges SOAP et REST. Le cinquième propose une discussion comparative et un retour d'expérience. Enfin, le sixième chapitre conclut sur les apports pédagogiques et ouvre quelques perspectives.
